% Français 9115, batch 9 — Gallica views 161–180.
% Pass: Opus 5 (claude-opus-5), 2026-09-15 — first pass, unchecked against
% the views by a human.
%
% What the twenty views hold. One thing, from the first view to the last:
% the copy, begun at view 144 (batch 8), of Euler's « Investigatio motuum
% quibus laminæ et virgæ elasticæ contremiscunt » (Acta Petropolitana for
% 1779, pars prior, as the copy's first leaf gives it). Same regular fair
% hand, written on both sides of each leaf, one parenthesised page number to
% a view, (120) at view 161 to (139) at view 180. All twenty views carry
% text; none is skipped. The copy runs on into batch 10.
%
%   v. 161      Case I (both ends free), end of the solution: $f = a/\omega$,
%               the time of one vibration and the sound, $\alpha, \beta,
%               \gamma, \delta$ from the two signs of $\sin\omega$, and the
%               equation of motion with $u = s/a$.
%   v. 162–163  Corollaries 1–4, §§ 16–19: no root in the first quadrant; one
%               near $3\zeta + 1^\circ$ ($\zeta$ the right angle, $e^{3\zeta}
%               \approx 111{,}31$), then $5\zeta, 7\zeta, 9\zeta \ldots$; the
%               sounds go as $9, 25, 49, 81$, the second an octave and a
%               tritone above the first (« b » against « cis »).
%   v. 164–166  Scholia 1–3, §§ 20–22: the general motion as a sum over the
%               roots; the fundamental has two nodes (« Tab. II. fig. 6 »,
%               curve $eLOMf$); the ordinates and tangents at $E$, $O$, $F$.
%   v. 168      « Evolutio casus II quo alter terminus liber relinquitur,
%               alter vero circa stylum mobilis figitur », Problema § 23.
%   v. 169, 167 Its solution to $\tan\omega = (e^\omega - e^{-\omega})/(e^\omega
%               + e^{-\omega})$ and the equation of motion.
%   v. 170–172  No root in the first quadrant; §§ 24–27, roots
%               $5\pi/4, 9\pi/4, 13\pi/4 \ldots$, sounds as $25, 81, 169,
%               289$; the fundamental against case I as $25 : 36$, « falsa
%               quinta » (g against a lower cis); Scholion 1, § 28, the general
%               motion.
%   v. 172–174  Scholion 2, § 29: the figure of the fundamental, ordinates
%               computed at $u = 0, \frac15, \ldots, \frac45$ with
%               $e^{2\omega} = 2576$; one node at $O$ (« Tab. II. fig. 7 »).
%   v. 175–177  « Evolutio casus III », one end free, the other « muro
%               firmiter infixus », Problema § 30: $\cos\omega = -2/(e^\omega +
%               e^{-\omega})$, $\alpha, \beta, \gamma, \delta$, time and sound;
%               Corollary I, § 31, the root of the second quadrant by series,
%               $\varphi = 23^\circ 29'$ then $16^\circ 58'$.
%   v. 178–179  Corollary 2, § 32: a trial-and-logarithm method, $\varphi =
%               12^\circ$ then $15^\circ$, giving $\varphi = 16^\circ 45'$,
%               $\omega = 106^\circ 45' \approx \frac35\pi$; the fundamental
%               two octaves and more below case I (« Tab. II fig. 8 »).
%               Corollary 3, § 33 begins.
%   v. 180      § 33: the further roots $5\zeta, 7\zeta, 9\zeta$, whose sounds
%               agree with case I; Scholion § 34, the general equation.
%
% The order of the leaves. The copy's pages are bound out of order at one
% place: view 166 is (125), view 167 is (128), and views 168–169 are (126)
% and (127). The text follows the page numbers — 166 runs into 168, 169 into
% 167, 167 into 170 — and the copyist says so herself at the foot of view 166:
% « (voyez p. 126 après 128) ». The views are transcribed in their Gallica
% order, as the skill requires, and the notes say where the sentence goes.
%
% The numbers on the leaves, and why no \folio{} is written. The ink series
% of the earlier batches runs on, at the top right of the even views: 80,
% 81, 82, 83, 84, 85, 86, 87, 88, 89 at views 162, 164, 166, 168, 170, 172,
% 174, 176, 178, 180 (79 at view 160, batch 8). The odd views carry no ink
% number. The copy's parenthesised page numbers, a third kind, are given at
% the head of each view: (120)–(125) at 161–166, (128) at 167, (126) at 168,
% (127) at 169, (129)–(139) at 170–180. That they report the printed pages of
% the Acta is batch 8's inference, not checked here. The oblique pencil
% strokes on the mounting sheets that batch 8 describes are there again and
% are not a figure. No pencilled foliation of the Bibliothèque was read, and
% no \folio{} is written. Every number above was read on its view.
%
% Notation. The copy's letters are kept, as in batch 8: $\alpha, \beta,
% \gamma, \delta$, $f$, $k$, $u = s/a$, « sin. » and « cos. » with their stops
% where written (set \text{sin.} and \cos.), « tang » set \operatorname{tang},
% $l$ for the logarithm. One looped letter serves both for the phase in
% $\sin(\zeta + t\ldots)$ and for the right angle (« signum anguli recti »,
% « character anguli recti »); it is set $\zeta$ in both uses. The
% copyist's underlining of single letters and of note names is set \emph{}.
% Signs and figures are as read; where they disagree with the calculation
% (views 161, 167, 169, 171, 172, 174, 176) the leaf stands and a note says
% so. Nothing is corrected.
%
% The hand, and one trap in it. The figure 3 of this hand is a short 3
% whose tail is easily taken for a 9's; the 9 has a long descending tail.
% Numbers were read with that distinction and checked where the calculation
% allows (e.g. $e^{\pi/4} = 2{,}1933$, $\frac35\pi = 108^\circ$, $m = 0{,}26180$).
%
% What could not be read. Nothing is marked \ill{}. Readings offered with
% doubt: $\Phi$ and « 1 » in the abbreviations of view 161; « sonus simul
% auditus » at 163; « Haec » and « debeat » at 164; « referenda » at 165;
% « deteguntur » and the last figure of « 128 » at 166; a sign at 169;
% « purum » at 173; « semisse » at 179.
%
% Nothing here was checked against any published transcription or edition
% of Euler's memoir.
\documentclass[11pt,a4paper]{article}
% The preamble shared by every transcription of the mathematicians' manuscripts in Gallica.
%
% It rests on one idea: **uncertainty compounds, it does not resolve.** A
% transcription that smooths over a doubtful word has destroyed the only thing
% it was made for — a reader must be able to tell, at every point, what was on
% the page and what was supplied. The apparatus macros below are therefore the
% critical apparatus, not decoration.
%
% The same commands are recognised by `scripts/render.mjs`, which produces the
% site's reading view, and by `scripts/tei.mjs`. A macro added here must be
% added there too, or rendering fails loudly — which is the wanted behaviour.
%
% Comments here are English, like the rest of the repository. The documents
% this preamble sets are French, and that is deliberate: see the skills.

\ProvidesPackage{germain}[2026/09/15 Transcriptions of mathematicians' manuscripts in Gallica]

\usepackage{amsmath,amssymb,amsthm}
\usepackage{mathrsfs}
% Kept from the parent preamble so the renderer's diagram support stays
% available; these manuscripts draw no commutative diagrams, and nothing here uses it.
\usepackage{tikz-cd}
\usepackage[english,french]{babel}
\usepackage{fontspec}

% --- Greek ------------------------------------------------------------------
% The text face stays fontspec's default, Latin Modern, which has no Greek:
% XeTeX drops every glyph it lacks with a line in the log and nothing on the
% page, and the Greek words the manuscripts quote vanished from the PDFs. So
% the Greek and Greek Extended blocks get a XeTeX character class of their own,
% and entering or leaving a run of it switches the family to CMU Serif — the
% same Computer Modern design, with polytonic Greek — and back. Only the
% family changes: a Greek word in \textit or \textbf keeps its shape and series.
% The transitions to and from every other class carry the tokens that class
% has towards an ordinary letter, so French spacing before « ; » or inside
% « » still holds after a Greek word. No grouping, so no brace or \endgroup
% can come out unbalanced; maths is left alone, since XeTeX inserts nothing
% there. Kept identical in `exercices.sty`.
\makeatletter
\newfontfamily\germainGreekfont{cmun}[
  Extension=.otf, NFSSFamily=germaingreek,
  UprightFont=*rm, ItalicFont=*ti, SlantedFont=*sl,
  BoldFont=*bx, BoldItalicFont=*bi, BoldSlantedFont=*bl]
\newXeTeXintercharclass\germain@greek
\def\germain@greekrange#1#2{%
  \@tempcnta=#1\relax
  \loop
    \XeTeXcharclass\@tempcnta=\germain@greek
    \ifnum\@tempcnta<#2\advance\@tempcnta\@ne\repeat}
\germain@greekrange{"0370}{"03FF}
\germain@greekrange{"1F00}{"1FFF}
\def\germain@greekprev{\rmdefault}
\def\germain@greekin{%
  \edef\germain@greekprev{\f@family}\fontfamily{germaingreek}\selectfont}
\def\germain@greekout{\fontfamily{\germain@greekprev}\selectfont}
\def\germain@greekjoin{%
  \XeTeXinterchartokenstate=\@ne
  \@tempcnta=\z@
  \loop
    \ifnum\@tempcnta=\germain@greek\else
      \XeTeXinterchartoks\germain@greek\@tempcnta=\expandafter{%
        \expandafter\germain@greekout\the\XeTeXinterchartoks\z@\@tempcnta}%
      \XeTeXinterchartoks\@tempcnta\germain@greek=\expandafter{%
        \the\XeTeXinterchartoks\@tempcnta\z@\germain@greekin}%
    \fi
    \ifnum\@tempcnta<4095\advance\@tempcnta\@ne\repeat}
% babel-french sets its own classes, and saves and restores the interchar
% state, each time a language is selected: join again after it does.
\addto\extrasfrench{\germain@greekjoin}
\addto\extrasenglish{\germain@greekjoin}
\AtBeginDocument{\germain@greekjoin}
\makeatother

\usepackage{geometry}
\geometry{a4paper,margin=25mm,marginparwidth=28mm}
\usepackage{marginnote}
\usepackage{xcolor}
\usepackage[normalem]{ulem}
\setlength{\emergencystretch}{3em}
\usepackage{hyperref}
\hypersetup{colorlinks=true,linkcolor=black,urlcolor=black,pdfborder={0 0 0}}

% --- Batch metadata ---------------------------------------------------------
% Taken as they stand by the rendering script, which shows them at the head of
% the reading view. They fix what the file claims to cover. The macro names are
% the parent project's, so that the scripts stay shared; \folder here means the
% volume's slug — `fr-9115` — and \pages counts Gallica views.

\newcommand{\folder}[1]{\gdef\grothFolder{#1}}
\newcommand{\batch}[1]{\gdef\grothBatch{#1}}
\newcommand{\pages}[2]{\gdef\grothFirst{#1}\gdef\grothLast{#2}}
\newcommand{\foldertitle}[1]{\gdef\grothFoldertitle{#1}}

% The holder's own shelfmark, printed as they write it: « Français 9115 ».
\newcommand{\shelfmark}[1]{\gdef\grothShelfmark{#1}}

% The dating, copied verbatim from the holder's catalogue — never inferred
% here. « XIXe siècle » is the BnF's; a pass that dates a leaf from its content
% says so in a \note{}, as its own inference.
\newcommand{\dating}[1]{\gdef\grothDating{#1}}

% The status notice, printed in the title block and repeated as a diagonal
% watermark on every page of the PDF. The manuscripts are in the public
% domain; what this line declares is not a right but a quality — an
% unverified first pass by a machine. The skills write the line; a file
% without it gets no watermark, so leaving it out is visible in the output.
\newcommand{\watermark}[1]{\gdef\grothWatermark{#1}}

\gdef\grothFolder{?}\gdef\grothBatch{}\gdef\grothFirst{?}\gdef\grothLast{?}%
\gdef\grothFoldertitle{}\gdef\grothShelfmark{}\gdef\grothDating{}\gdef\grothWatermark{}

% --- The résumé -------------------------------------------------------------
\newenvironment{resume}
  {\small\setlength{\parskip}{0.45em}}
  {\par\medskip\hrule\medskip}

% --- Keywords ---------------------------------------------------------------
% In English, closing the modernised reading's résumé: the modern vocabulary
% under which to search for what the leaves construct. The manifest extracts
% them to tag the volume — this line is the single source of the tags.
\newcommand{\keywords}[1]{\par\smallskip\noindent
  {\footnotesize\textsc{Keywords} --- \textit{#1}}\par}

% --- The critical apparatus -------------------------------------------------

% The Gallica view number, in the margin. This is the anchor that lets a
% disagreement over a formula be settled by looking at *one* image rather than
% a volume of 750. The site uses it to turn the facsimile as the transcription
% is scrolled. A view is one scanned image: usually one side of a leaf.
\newcommand{\page}[1]{%
  \marginnote{\footnotesize\color{gray!70}\textsf{v.\,#1}}%
  \label{page:#1}\ignorespaces}

% The BnF's pencilled foliation, where the leaf carries it and a pass can read
% it: « 198r ». This is what the literature cites — « ff. 198r–208v » — and
% the one line that turns a citation into a view. Recorded, never inferred:
% a folio a pass cannot read is not written.
\newcommand{\folio}[1]{%
  \marginnote{\footnotesize\color{gray!50}\textsf{f.\,#1}}\ignorespaces}

% The modernised reading's coarser counterpart to \page{}: which views a
% section is drawn from.
\newcommand{\pagerange}[2]{%
  \marginnote{\footnotesize\color{gray!70}\textsf{v.\,#1--#2}}%
  \ignorespaces}

% Illegible: never guessed.
\newcommand{\ill}{\textcolor{red!60!black}{[\,\ldots\,]}}

% A doubtful reading: the word is given, and flagged as uncertain.
\newcommand{\uncertain}[1]{\uline{#1}}

% An editorial addition — a missing letter, an implied word.
\newcommand{\add}[1]{\textcolor{blue!50!black}{[#1]}}

% Struck out by the writer. What was crossed out is part of the document.
\newcommand{\struck}[1]{\sout{#1}}

% The transcriber's note, on the state of the page or the mathematics.
\newcommand{\note}[1]{\par\smallskip{\small\color{gray!60!black}\textit{[#1]}}\par\smallskip}

% A marginal note of hers, distinct from the body of the text.
\newcommand{\marginal}[1]{\par\smallskip\noindent\hspace{1em}%
  {\small\color{gray!60!black}#1}\par\smallskip}

% --- Title ------------------------------------------------------------------
% The title block is French, like the documents themselves.

\AddToHook{shipout/background}{%
  \ifx\grothWatermark\empty\else
    \begin{tikzpicture}[remember picture, overlay]
      \node[rotate=52, scale=3.4, align=center, text=gray!18,
            font=\sffamily\bfseries]
        at (current page.center) {\grothWatermark};
    \end{tikzpicture}%
  \fi}

\AtBeginDocument{%
  \begin{center}
    {\large\bfseries Manuscrits de mathématiciens (Gallica) ---
      \ifx\grothShelfmark\empty\grothFolder\else\grothShelfmark\fi}\\[2pt]
    {\itshape\grothFoldertitle}\\[4pt]
    {\small vues \grothFirst--\grothLast
      \ifx\grothBatch\empty\else\ (lot \grothBatch)\fi}\\[2pt]
    \ifx\grothDating\empty\else
      {\small\color{gray!60!black}Datation du catalogue : \grothDating}\\[2pt]
    \fi
    \ifx\grothWatermark\empty\else
      {\small\bfseries\color{red!45!gray}\grothWatermark}\\[2pt]
    \fi
    {\footnotesize\color{gray!60!black}Transcription établie d'après les images IIIF de Gallica.
     Source gallica.bnf.fr / Bibliothèque nationale de France.\\
     Les lectures incertaines sont soulignées, les illisibles notées [\,\ldots\,],
     les ajouts entre crochets ; \textsf{v.} numérote les vues de Gallica,
     \textsf{f.} la foliotation au crayon de la BnF.}
  \end{center}
  \vspace{1em}}

\endinput


\folder{fr-9115}
\shelfmark{Français 9115}
\batch{9}
\pages{161}{180}
\dating{1801-1900}
\watermark{Lecture automatique\\non vérifiée}
\foldertitle{Papiers de Sophie GERMAIN. — Recueil de dissertations et problèmes mathématiques et physiques.}

\begin{document}

\note{Numérotation des feuillets. La série à l'encre des lots précédents se
poursuit en haut à droite des vues paires : 80, 81, 82, 83, 84, 85, 86, 87, 88,
89 aux vues 162, 164, 166, 168, 170, 172, 174, 176, 178, 180. Les vues impaires
ne portent pas de nombre à l'encre. La copie latine porte en tête de chaque vue
un nombre entre parenthèses : (120) à (125) aux vues 161 à 166, (128) à la vue
167, (126) à la vue 168, (127) à la vue 169, (129) à (139) aux vues 170 à 180.
Les pages (126) à (128) sont donc reliées hors de leur ordre, ce que la copie
signale elle-même au bas de la vue 166. Tous ces nombres ont été lus sur la
vue ; aucun n'est déduit. Aucune foliotation de la Bibliothèque n'a été lue sur
ces vues, et aucun « f. » n'est donc porté ici.}

\page{161}
(120)

$f = \frac{a}{\omega}$, hincque $k = \frac{a^4}{bcc\omega^4}$; tum vero erit
$\sqrt{\frac{2g}{k}} = \frac{c\omega\omega\sqrt{2bg}}{aa}$;
quocirca tempus unius vibrationis erit $= \frac{\pi aa}{\omega\omega c\sqrt{2gb}}$,
hincque ipse sonus a virga editus $= \frac{\omega\omega c\sqrt{2gb}}{\pi aa}$.
Denique cum sit $\cos\omega = \frac{2}{e^\omega + e^{-\omega}}$, erit
$\text{sin}\,\omega = \frac{\pm(e^\omega - e^{-\omega})}{e^\omega + e^{-\omega}}$;
ex priore autem valore, quo sinus $\omega$ positivus, erit
\[
\frac{\alpha}{\beta} = \frac{1 - e^{-2\omega} - e^\omega + e^{-\omega}}{-1 + e^{2\omega} - e^\omega + e^{-\omega}},
\quad\text{hincque}
\]
\[
\frac{\alpha}{\beta}e^\omega = \frac{e^\omega - e^{-\omega} - e^{2\omega} + 1}{-e^\omega + e^{-\omega} + e^{2\omega} - 1} = -1,
\]
ita ut sit $\frac{\alpha}{\beta} = \frac{-1}{e^\omega} = -e^{-\omega}$; altero autem
casu quo sinus $\omega$ est negativus, reperitur
$\frac{\alpha}{\beta} = +\frac{1}{e^\omega} = +e^{-\omega}$; quare si sumamus $\alpha = 1$
erit $\beta = \mp e^\omega$, hincque $\gamma = 1 \pm e^\omega$ et $\delta = 1 \mp e^\omega$,
ubi signa superiora valent, si $\text{sin}\,\omega$ est positivus, inferiora si negativus.
His igitur valoribus inventis aequatio pro motu virgae perfecte est determinata,
quae quo simplicius repraesentetur, notetur, ob $f = \frac{a}{\uncertain{\Phi}}$ esse
$\frac{s}{f} = \frac{s}{\omega}$; unde si ponatur brevitatis gratia $\frac{\uncertain{1}}{a} = \omega$,
ut sit $\frac{s}{f} = u\omega$, hoc valore posito aequatio nostra pro motu erit:
\note{transcrit tel qu'on le lit sur la vue : le dénominateur de $f$ a la forme
d'un $\Phi$ (en tête de page, $f = \frac{a}{\omega}$), le second membre de
$\frac{s}{f}$ est écrit $\frac{s}{\omega}$, et la définition abrégée
$\frac{1}{a} = \omega$, alors que la suite ($\frac{s}{f} = u\omega$) demande
$\frac{s\omega}{a}$ et $\frac{s}{a} = u$. Le numérateur « 1 » peut être un
$s$ de grande taille.}
\[
y = C\text{sin}\,\Bigl(\zeta + t\,\frac{\omega\omega c\sqrt{2gb}}{aa}\Bigr)\bigl(e^{u\omega} \mp e^{(1-u)\omega} + (1 \pm e^\omega)\text{sin.}\,u\omega + (1 \mp e^\omega)\cos.u\omega\bigr)
\]
ex qua ad quodvis tempus $t$ status virgae elasticae cognoscitur, si modo valor
literae $\omega$ rite fuerit definitus.

\page{162}
(121)

§.~16. Cum igitur $\omega$ denotet arcum circuli, cujus radius est $= 1$, in primo
quadrante haud difficulter perspicitur, post casum $\omega = 0$ nullum alium angulum
satisfacere; semper enim erit $\cos\omega < \frac{2}{e^\omega + e^{-\omega}}$, quod ita
ostendi potest. Cum sit per series
\[
\cos\omega = 1 - \frac{\omega\omega}{1.2} + \frac{\omega^4}{1.2.3.4} - \frac{\omega^6}{1.2.3.4.5.6} + \&c. \quad\text{et}
\]
\[
e^\omega + e^{-\omega} = 2\Bigl(1 + \frac{\omega\omega}{1.2} + \frac{\omega^4}{1.2.3.4} + \frac{\omega^6}{1.2.3.4.5.6} + \&c\Bigr) \quad\text{erit}
\]
\[
\cos\omega\,(e^\omega + e^{-\omega}) = 2\Bigl(1 - \frac{\omega^4}{6}\Bigr)
\]
qui valor manifesto minor est quam $2$, nisi valor ipsius $\omega$ angulum rectum
superet.

\subsection*{Corollarium 2.}

§.~17. Deinde manifestum est, neque in secundo neque in tertio quadrante reperiri
valorem ipsius $\omega$, quia in his quadrantibus cosinus sunt negativi, formula
autem $\frac{2}{e^\omega + e^{-\omega}}$ semper est positiva. At in quarto quadrante,
ubi cosinus iterum fiunt positivi, dabitur valor non multum tres angulos rectos
superans. Sit igitur $\zeta$ signum anguli recti, sive $\zeta = \frac{\pi}{2}$, ac
ponatur $\omega = 3\zeta + \varphi$, fietque $\cos\omega = \text{sin.}\,\varphi$, unde fieri debet
\[
\text{sin.}\,\varphi = \frac{2}{e^{3\zeta+\varphi} + e^{-3\zeta-\varphi}};
\]
ubi facile perspicitur, angulum $\varphi$ valde esse parvum, quia numerus
$e^{3\zeta}$, ob $3\zeta = 4{,}71239$ et $e = 2{,}71828$, est satis magnus, scilicet
proxime $= 111{,}31$, pro quo scribamus \emph{n}. Jam quia $\text{sin}\,\varphi = \varphi$,
proxime, $e^\varphi = 1 + \varphi$ et $e^{-\varphi} = 1 = 1 - \varphi$ nostra aequatio erit
\[
\varphi = \frac{2n}{nn(1+\varphi) + 1 - \varphi}, \quad\text{sive}\quad \varphi = \frac{2n}{nn+1} \ \text{proxime};
\]
\note{« $e^{-\varphi} = 1 = 1 - \varphi$ » est écrit ainsi sur le feuillet. La lettre
rendue ici par $\zeta$ est un signe bouclé en forme de $\zeta$, pris pour
« signum anguli recti ».}

\page{163}
(122)

sumi igitur circiter poterit $\varphi = \frac{2}{n} = \frac{1}{56}$; unde patet angulum
$\varphi$ vix unum gradum superare, ita ut sit $\omega = 3\zeta + 1^\circ$ circiter

\subsection{Corollarium 3.}

§.~18. Progrediamur ad quintum quadrantem, et manifestum est hic iterum dari
valorem tantillo minorem quam $5\zeta$; defectus enim aliquot minuta prima non
excedet. Porro vero sextus et septimus vacui manebunt; in octavo autem reperietur
$\omega = 7\zeta$, tam prope ut excessus sentiri nequeat; sicque valores ulteriores
erunt continuo exactius $\omega = 9\zeta$, $\omega = 11\zeta$, \&c.

\subsection{Corollarium 4}

§.~19. Quod si ergo in primo valore exiguum discrimen unius gradus negligamus,
omnes valores anguli $\omega$ erunt $3\zeta$, $5\zeta$, $7\zeta$, $9\zeta$, $11\zeta$,
\&c qui secundum numeros impares in infinitum progrediuntur, unde nostra virga
infinitos sonos simplices edere poterit, qui sequentibus numeris exprimentur:
\[
\frac{9\zeta c\sqrt{2gb}}{2aa};\quad \frac{25\zeta c\sqrt{2gb}}{2aa};\quad
\frac{49\zeta c\sqrt{2gb}}{2aa};\quad \frac{81\zeta c\sqrt{2gb}}{2aa};\quad \&c.
\]
qui ergo soni secundum numeros, $9$. $25$, $49$, $81$ \&c. progrediuntur, quorum
primus pro fundamentali haberi potest; proximus vero ad hunc rationem tenebit ut
$25 : 9$, quod intervallum complectitur unam octavam cum tritono. Unde si sonus
fundamentalis fuerit b, sequens futurus erit \emph{cis}, qui ergo
\uncertain{sonus simul auditus} valde ingratam dissonantiam referat.
\note{« cis » est souligné sur le feuillet.}

\page{164}
(123)

\subsection{Scholion 1.}

§.~20. Quodsi igitur pro singulis istis valoribus anguli $\omega$ formentur valores
ipsius $y$, eorum quotcunque invicem conjuncti exhibebunt motus, quos nostra virga
recipere poterit. Si hoc modo omnes infiniti valores ipsius $\omega$ invicem
conjungantur, ac pro quolibet literis $C$ et $\zeta$ generatim quicunque alii
valores tribuantur, aequatio obtinebitur generalis, quae omnes plane motus, qui in
virga locum habere possunt, in se complectetur. \uncertain{Haec} igitur designemus, sequentem
habebit formam:
\begin{align}
y ={}& C\text{sin}\,\Bigl(\zeta + t\,\frac{\omega\omega c\sqrt{2gb}}{aa}\Bigr)\bigl\{e^{u\omega} \mp e^{(1-u)\omega} + (1 \pm e^\omega)\text{sin.}\,u\omega + (1 \mp e^\omega)\cos.u\omega\bigr\}\\
&+ C'\text{sin}\,\Bigl(\zeta + t\,\frac{\omega'\omega'c\sqrt{2gb}}{aa}\Bigr)\bigl\{e^{u\omega'} \mp e^{(1-u)\omega'} + (1 \pm e^{\omega'})\text{sin.}\,u\omega' + (1 \mp e^{\omega'})\cos.u\omega'\bigr\}\\
&+ C''\text{sin}\,\Bigl(\zeta + t.\frac{\omega''\omega''c\sqrt{2gb}}{aa}\Bigr)\bigl\{e^{u\omega''} \mp e^{(1-u)\omega''} + (1 \pm e^{\omega''})\text{sin.}\,u\omega'' + (1 \mp e^{\omega''})\cos.u\omega''\bigr\}\\
&\&c \qquad\qquad \&c
\end{align}
ubi litera $u$ denotet fractionem $\frac{s}{a}$, et signa superiora valeant si
angulorum $\omega$, $\omega'$, $\omega''$ \&c sinus fuerint positivi, inferiora vero
si fuerint negativi

\subsection{Scholion 2}

§.~21. Quod si sonum fundamentalem, quo est proxime $\omega = 3\zeta$, et qui
continet motum simplicissimum, quo virga contremiscere potest, attentius
consideremus, facile colligere licet, curvam, quam virga inter vibrandum induit,
axem ad minimum in duobus punctis secare \uncertain{debeat}, ita ut quasi duos nodos formet.
Si enim axem
\note{la fin de « debeat » est reprise à la plume par-dessus d'autres lettres ;
le « ut » de « ita ut » est de même surchargé.}

\page{165}
(124)

nusquam secaret, dum singula ejus puncta ab axe recedunt, eodem motu continuo
ulterius recedere deberent, quia virga a nullis plane viribus coercetur; quod inde
etiam perspicuum est quod hoc casu centrum gravitatis virgae immotum esse debeat.
Si porro virga in motu suo unicum nodum formaret, circa quem quasi gyraretur,
motum semel conceptum gyratorium perpetuo conservare deberet. Hinc igitur patet,
virgam inter vibrandum ejusmodi formam $eLOMf$ esse habituram,
\marginal{Tab. II. fig. 6.}
quae situm naturalem, seu axem $EF$ in duobus punctis $L$, $M$ secat. Hoc idem vero
etiam nostra formula declarat: quia enim angulus $\omega$ hic subito tres angulos
rectos superat, anguli $u\omega$, sive $\frac{\omega s}{a}$, dum quantitas $s$ usque
ad $a$ augetur, angulos \uncertain{referenda} $0$ usque ad $3\zeta$ continuo
ascendentes, quorum sinus et cosinus interea bis contraria signa recipiunt, unde
duobus casibus contingere potest, ut applicata $y$ evanescat. Eodem modo intelligere
licet, pro secundo ipsius $\omega$ valore $= 5\zeta$ curvam virgae tres nodos habere
debere, pro sequente, $\omega = 7\zeta$, quatuor, et ita porro. Singuli autem isti
nodi sive intersectiones cum axe $EF$ pro quovis valore ipsius $\omega$ ex hac
aequatione elici poterunt:
\[
e^{u\omega} \mp e^{(1-u)\omega} + (1 \pm e^\omega)\text{sin.}\,\omega u + (1 \mp e^\omega)\cos.\omega u = 0
\]
quippe ex qua valores literae $u = \frac{s}{a}$ ipsos nodos declarabunt. Scilicet
literae $u$, incipiendo a $0$, continuo majores tribuantur valores usque ad $1$, et
casus notentur, quibus ista formula evanescit; si enim quispiam valor jam evadat
valde parvus per regulam approximationis
\note{« nusquam », au début de la page, est écrit par-dessus un autre mot ;
« plane », « patet », « aequatione » et « poterunt » portent des lettres
reprises à la plume. « Tab. II. fig. 6. » est dans la marge de gauche, en
regard de la ligne qui nomme la figure $eLOMf$ ; aucune planche n'est sur ces
vues.}

\page{166}
(125)

veri ejus valores facile \uncertain{deteguntur}.

\subsection{Scholion 3}

§.~22. Quo haec clarius perspiciantur, casum primum, quo $\omega = 3\zeta$ circiter,
ideoque ejus sinus negativus, accuratius perpendamus. Erit igitur, primo factore,
constante seu a tempore pendente, omisso:
\[
y = \cdots\ e^{u\omega} + e^{\omega(1-u)} + (1 - e^\omega)\text{sin}\,u\omega + (1 + e^\omega)\cos u\omega,
\]
cujus valores pro tribus casibus $u = 0$, $u = 1$, $u = \frac{1}{2}$ definiamus, ut
applicatas non solum pro utroque termino $E$ et $F$, sed etiam pro puncto medio
\emph{O}, scilicet $Ee$, $Ff$, et $Oo$, obtineamus. Primo igitur, posito $u = 0$ erit
$Ee = 2(1 + e^\omega)$; posito autem $u = 1$, prodit applicata
\[
Ff = e^\omega + 1 + (1 - e^\omega)\text{sin}\,\omega + (1 + e^\omega)\cos.\omega,
\]
qui valor, ob $\text{sin.}\,\omega = -1$ et $\cos\omega = \frac{2}{e^\omega + e^{-\omega}}$,
abit in hunc: $Ff = 2(1 + e^\omega)$, sicque hae duae applicatae $Ee$ et $Ff$ inter se
erunt aequales; at pro puncto medio $O$, ubi fit $u = \frac{1}{2}$, hincque
$\text{sin.}\,\frac{1}{2}\omega = \text{sin.}\,\frac{3}{2}\zeta = \frac{1}{\sqrt{2}}$ et
$\cos\frac{\omega}{2} = -\frac{1}{\sqrt{2}}$, erit applicata
\[
Oo = e^{\frac{\omega}{2}} + e^{\frac{\omega}{2}} + \frac{(1 - e^\omega)}{\sqrt{2}} - \frac{(1 + e^\omega)}{\sqrt{2}} = 2e^{\frac{\omega}{2}} - \sqrt{2}.e^\omega,
\]
qui valor, ob $e^\omega = 111$ et $e^{\frac{\omega}{2}} = 10\frac{1}{2}$, abit in
$21 - 111.\sqrt{2}$; unde patet, hanc applicatam esse negativam, prorsus uti figura
refert. Possumus etiam simili modo positionem tangentium pro his locis exhibere ex
formula
\[
\Bigl(\frac{dy}{ds}\Bigr) = \cdots\ e^{u\omega} - e^{\omega(1-u)} + (1 - e^\omega)\cos.u\omega - (1 + e^\omega)\text{sin.}\,u\omega,
\]
quae formula, posito $u = 0$, pro termino $E$ praebet:
\[
\Bigl(\frac{dy}{ds}\Bigr) = 1 - e^\omega + 1 - e^\omega = 2(1 - e^\omega) = -2(e^\omega - 1);
\]
\marginal{(voyez p.~126 après 12\uncertain{8})}
\note{« \emph{O} » est souligné. La mention entre parenthèses, en français, est
écrite à la fin de la dernière ligne, d'une plume plus fine ; le dernier chiffre de
« 128 » a la même forme que le 6 de « 126 » et n'est pas sûr. La page qui suit dans
la reliure (vue 167) porte « (128) » et continue le cas II ; la suite du présent
calcul, « tum vero, posito $u = 1$ », est à la page « (126) », vue 168.}

\page{167}
(128)

ideoque $\frac{\alpha}{\beta}e^{2\omega} = -1$, unde fit $\frac{\alpha}{\beta} = -e^{-2\omega}$.
Posteriore vero casu colligitur
\[
\frac{\alpha}{\beta} = \frac{-e^{-\omega}\sqrt{2}(e^{2\omega} + e^{-2\omega}) + 2e^{-\omega}}{e^\omega\sqrt{2}(e^{2\omega} + e^{-2\omega}) - 2e^\omega} = \frac{-1}{e^{2\omega}} = -e^{-\omega}.
\]
Utroque ergo casu, sive sinus et cosinus anguli \emph{$\omega$} sint ambo positivi,
sive ambo negativi, pro fractione $\frac{\alpha}{\beta}$ idem valor obtinetur. Quare
si ponatur $\alpha = 1$, erit $\beta = -e^{2\omega}$, hincque $\gamma = 1 + e^{2\omega}$
et $\delta = 1 - e^{2\omega}$; quibus valoribus inventis aequatio pro motu virgae, si
iterum loco $\frac{s}{a}$ scribamus \emph{u}, erit haec
\[
y = C\text{sin}\,\Bigl(\zeta + t\sqrt{\frac{2g}{k}}\Bigr)\bigl\{e^{u\omega} - e^{\omega(2-u)} + (1 + e^{2\omega})\text{sin.}\,u\omega + (1 - e^\omega)\cos.u\omega\bigr\}.
\]
\note{au dernier membre de la fraction, « $-e^{-\omega}$ », et au dernier terme de
l'équation, « $(1 - e^\omega)$ », l'exposant est écrit sans le 2 que portent les
lignes précédentes ($-e^{-2\omega}$, $\delta = 1 - e^{2\omega}$) ; transcrit tel quel.}

Quicunque autem valores pro angulo $\omega$ ex aequatione
\[
\operatorname{tang}\omega = \frac{e^\omega - e^{-\omega}}{e^\omega + e^{-\omega}}
\]
eruentur, ex singulis habetur
\[
f = \frac{a}{\omega}, \ \text{hincque}\ f^4 = \frac{a^4}{\omega^4} = bcck, \ \text{unde fit}
\]
\[
k = \frac{a^4}{bcc\omega^4} \ \text{et}\ \sqrt{\frac{2g}{k}} = \frac{\omega\omega c\sqrt{2gb}}{aa},
\]
hincque porro ut ante tempus unius oscillationis erit:
\[
\pi\sqrt{\frac{k}{2g}} = \frac{\pi aa}{\omega\omega c\sqrt{2gb}},
\]
et sonus a virga elastica editus $= \frac{\omega\omega c\sqrt{2gb}}{\pi aa}$. Totum ergo
negotium huc est reductum, ut ex aequatione
\[
\operatorname{tang}.\omega = \frac{e^\omega - e^{-\omega}}{e^\omega + e^{-\omega}}
\]
omnes valores anguli $\omega$ eruantur; ubi quidem statim liquet, valorem
$\omega = 0$ satisfacere, unde autem nullus motus sequitur; quare pro reliquis
valoribus singulos quadrantes percurramus. Pro primo igitur quadrante per series
habebimus:

\page{168}
(126)

tum vero, posito $u = 1$, pro termino $F$ erit
\[
\Bigl(\frac{dy}{ds}\Bigr) = e^\omega - 1 + (1 - e^\omega)\cos.\omega - (1 + e^\omega)\text{sin.}\,\omega,
\]
qui valor, ob $\cos.\omega = -1$ et $\text{sin.}\,\omega = \frac{2}{e^\omega}$ rejecto termino
$e^{-\omega}$, ut pote minimo, reducitur ad $\bigl(\frac{dy}{ds}\bigr) = 2(e^\omega - 1)$;
unde patet angulos $EeL$ et $FfM$ esse inter se aequales. Pro puncto autem medio $O$,
ubi $u = \frac{1}{2}$, prodit
\[
\Bigl(\frac{dy}{ds}\Bigr) = (1 - e^\omega)\cos\tfrac{1}{2}\omega - (1 + e^\omega)\text{sin.}\,\tfrac{1}{2}\omega = -\sqrt{2},
\]
qui valor, si calculus accuratius institueretur, prodiret $= 0$, ita ut tangens in
puncto \emph{o} axi sit parallela. Simili prorsus modo etiam sequentes casus, ubi
$\omega = 5\zeta$ vel $7\zeta$ vel $9\zeta$ expendere licebit.
\note{dans la première formule, le $\omega$ de « $\cos.\omega$ » est repris à la plume.}

\subsection{Evolutio casus II quo alter terminus liber relinquitur, alter vero, circa stylum mobilis figitur}

\subsection{Problema}

§.~23. Si virga elastica in termino $E$ fuerit libera, in altero vero $F$ stylo
affixa, circa quem tamen libere moveri possit, investigare omnes vibrationes
regulares, quibus ea contremiscere potest.

\subsection{Solutio}

Quia ergo pro termino $E$, ubi $s = 0$, ut ante est
$\bigl(\frac{ddy}{ds^2}\bigr) = 0$ et $\bigl(\frac{d^3y}{ds^3}\bigr) = 0$, erit etiam ut
ante $\alpha + \beta - \delta = 0$ et $\alpha - \beta - \gamma = 0$, unde fit
$\gamma = \alpha - \beta$ et $\delta = \alpha + \beta$. Pro altero termino simpliciter
fixo, ubi $s = a$, posito

\page{169}
(127)

iterum $\frac{a}{f} = \omega$, primo debet esse $y = 0$, tum vero etiam
$\bigl(\frac{ddy}{ds^2}\bigr) = 0$. Prior conditio dat
\[
\alpha e^\omega + \beta e^{-\omega} + \gamma\text{sin.}\,\omega + \delta\cos.\omega = 0
\]
posterior vero
\[
\alpha e^\omega + \beta e^{-\omega} - \gamma\text{sin.}\,\omega - \delta\cos.\omega = 0
\]
quae aequationes, loco $\gamma$ et $\delta$ substitutis valoribus, abeunt in sequentes:
\[
\alpha(e^\omega + \text{sin.}\,\omega + \cos.\omega) + \beta(e^{-\omega} - \text{sin.}\,\omega + \cos.\omega) = 0 \quad\text{et}
\]
\[
\alpha(e^\omega - \text{sin.}\,\omega - \cos.\omega) + \beta(e^{-\omega} + \text{sin.}\,\omega - \cos.\omega) = 0;
\]
unde geminus valor oritur
\[
\frac{\alpha}{\beta} = -\frac{(e^{-\omega} - \text{sin}\,\omega + \cos.\omega)}{e^\omega + \text{sin}\,\omega + \cos.\omega} = -\frac{(e^{-\omega} + \text{sin.}\,\omega - \cos.\omega)}{e^\omega - \text{sin.}\,\omega - \cos.\omega}.
\]
Ponatur iterum $\text{sin.}\,\omega + \cos.\omega = p$ et $\text{sin}\,\omega - \cos.\omega = q$, ut sit
\[
\frac{-e^{-\omega} + q}{e^\omega + p} = \frac{-e^{-\omega} - q}{e^\omega \uncertain{+} p},
\]
unde colligitur haec aequatio:
\[
-1 + pe^{-\omega} + qe^\omega - pq = -1 - qe^\omega - pe^{-\omega} - pq
\]
sive $pe^{-\omega} + qe^\omega = 0$, unde concluditur
$\operatorname{tang}.\omega = \frac{e^\omega - e^{-\omega}}{e^\omega + e^{-\omega}}$
quocirca habebimus, vel
\[
\text{sin}\,\omega = \frac{e^\omega - e^{-\omega}}{\sqrt{2}(e^{2\omega} + e^{-2\omega})} \quad\text{et}\quad \cos.\omega = \frac{e^\omega + e^{-\omega}}{\sqrt{2}(e^{2\omega} + e^{-2\omega})}, \ \text{vel}
\]
\[
\text{sin.}\,\omega = \frac{-e^\omega + e^{-\omega}}{\sqrt{2}(e^{2\omega} + e^{-2\omega})} \quad\text{et}\quad \cos.\omega = \frac{-e^\omega - e^{-\omega}}{\sqrt{2}(e^{2\omega} + e^{-2\omega})}.
\]
ex prioribus valoribus colligitur
\[
\frac{\alpha}{\beta} = \frac{-e^{-\omega}\sqrt{2}(e^{2\omega} + e^{-2\omega}) - 2e^{-\omega}}{e^\omega\sqrt{2}(e^{2\omega} + e^{-2\omega}) + 2e^\omega}
\]
\note{au second dénominateur de la proportion, le signe devant $p$ ressemble à un
« + », comme au premier, alors que la ligne précédente donne $e^\omega - \text{sin.}\,\omega
- \cos.\omega$. Le radical des valeurs de $\text{sin.}\,\omega$ et $\cos.\omega$ est
transcrit $\sqrt{2}(\ldots)$ : on ne voit pas sur la vue s'il couvre la
parenthèse.}

\page{170}
(129)

\[
\text{sin}\,\omega\,(e^\omega + e^{-\omega}) = 2\omega\Bigl(1 + \frac{1}{3}\omega\omega - \frac{1}{30}\omega^4\Bigr) \quad\text{et}
\]
\[
\cos.\omega\,(e^\omega - e^{-\omega}) = 2\omega\Bigl(1 - \frac{1}{3}\omega\omega - \frac{1}{30}\omega^4\Bigr);
\]
unde patet, priorem formulam per totum primum quadrantem majorem esse quam
posteriorem, ita ut in hoc quadrante nullus reperiatur valor pro angulo $\omega$.
In secundo autem quadrante, ubi omnes tangentes sunt negativi, nullus iterum dari
potest, neque etiam in quarto, sexto, octavo et omnibus paribus. Reliquos quadrantes
in corollariis percurramus.
\note{« unde » et « esse », en tête de deux lignes, sont d'une écriture plus
petite ; le « In » de « In secundo » est écrit sur une autre lettre.}

\subsection{Corollarium 1.}

§.~24. Consideremus igitur tertium quadrantem, ubi cum sit $\omega$ majus quam
$\pi$, formula $\frac{e^\omega - e^{-\omega}}{e^\omega + e^{-\omega}}$ parum ab unitate
deficiet, unde angulus $\omega$ aliquando minor erit quam $\pi + 45^\circ$. Hinc
sumto iterum $\zeta$ pro signo anguli recti statuamus
$\omega = \pi + \frac{1}{2}\zeta - \varphi$, eritque
\[
\operatorname{tang}\omega = \operatorname{tang}\bigl(\tfrac{1}{2}\zeta - \varphi\bigr) = \frac{1 - \operatorname{tang}\varphi}{1 + \operatorname{tang}.\varphi}.
\]
Cum igitur sit
\[
\frac{1 - \operatorname{tang}\varphi}{1 + \operatorname{tang}\varphi} = \frac{e^\omega - e^{-\omega}}{e^\omega + e^{-\omega}} = \frac{1 - e^{-2\omega}}{1 + e^{-2\omega}},
\]
hinc, manifesto est $\operatorname{tang}\varphi = e^{-2\omega} = \frac{1}{e^{2\pi + \zeta - 2\varphi}}$,
in quo exponente angulum exiguum \emph{$\varphi$} negligere licet, ita ut subducto
calculo reperiatur $\operatorname{tang}\varphi = \frac{1}{2576}$; sicque angulus
$\varphi$ vix unum minutum superat, in quod tuto negligi potest ita ut primus valor
sit $\omega = \pi + \frac{1}{2}\zeta = 225^\circ$, cujus tam sinus quam cosinus est
$= -\frac{1}{\sqrt{2}}$.

\page{171}
(130)

\subsection{Corollarium 2}

§.~25. Pergamus igitur ad quartum quadrantem, ubi angulus $\omega$ multo minus
discrepabit a $2\pi + 45^\circ$. Hic autem tam sinus quam cosinus erit
$= +\frac{1}{\sqrt{2}}$. Simili modo ex septimo quadrante nanciscimur
$\omega = 3\pi + \frac{1}{2}\zeta$, tam sinu quam cosinu existente
$= -\frac{1}{\sqrt{2}}$; nonus vero quadrans suppeditat
$\omega = 4\pi + \frac{1}{2}\zeta$; undecimus $\omega = 5\pi + \frac{1}{2}\zeta$, et ita
porro.
\note{« quartum » est bien ce que porte le feuillet ; l'angle voisin de
$2\pi + 45^\circ$ appartient au cinquième quadrant.}

\subsection{Corollarium 3.}

§.~26. Cum igitur sit $\zeta = \frac{\pi}{4}$, omnes valores pro angulo
\emph{$\omega$} hactenus inventi sequenti modo progrediuntur:
\[
1^{\text{dus}}\ \omega = \frac{5\pi}{4};\quad 2^{\text{dus}}\ \omega = \frac{9\pi}{4};\quad
3^{\text{dus}}\ \omega = \frac{13\pi}{4};\quad 4^{\text{dus}}\ \omega = \frac{17\pi}{4};\ \&c.
\]
unde omnes soni simplices, quos ista virga edere potest, sequentibus numeris
exprimentur:
\[
\frac{25\pi c\sqrt{2gb}}{16aa};\quad \frac{81\pi c\sqrt{2gb}}{16aa};\quad
\frac{169\pi c\sqrt{2gb}}{16aa};\quad \frac{289\pi c\sqrt{2gb}}{16aa};
\]
quorum primus fundamentalis censetur; unde si secundus simul exaudiatur, erit primus
ad secundum ut $25$ ad $81$ hoc est ut $1 : 3\frac{6}{25}$, seu proxime ut
$1 : 3\frac{1}{4}$. Ergo si sonus fundamentalis fuerit $C$, sequens erit \emph{gis},
quem tamen sonum uno commate superabit, sicque harmonia parum grata existet
\note{« $\zeta = \frac{\pi}{4}$ » est écrit ainsi, alors que $\zeta$ désigne depuis
le §~16 l'angle droit, $\frac{\pi}{2}$.}

\subsection{Corollarium 4}

§.~27. Cum hic sonus fundamentalis, seu gravissimus, quem virga

\page{172}
(131)

haec edere potest sit $= \frac{25\pi c\sqrt{2gb}}{16aa}$, casu autem primo, quo
uterque terminus erat liber, sonus fundamentalis repertus fuerit
$\frac{9\zeta c\sqrt{2gb}}{2aa} = \frac{9\pi c\sqrt{2gb}}{4aa}$, ille se habebit ad
hunc ut $25 : 36$, ita ut casu primo eadem virga sonum edat fere una quinta
acutiorem. Scilicet si sonus casu primo editus fuerit \emph{g}, tum sonus secundo
casu editus erit gravior \emph{cis}, quod intervallum a musicis falsa quinta
appellatur. Hic igitur utique notari meretur, quod si virga utrinque libera edat
sonum \emph{g} tum eadem virga altero termino simpliciter fixa subito editura sit
sonum graviorem \emph{cis}, id quod experientia facile comprobari potest
\note{« cis » est souligné les deux fois ; « g » ne l'est pas sur le feuillet et
n'est mis en italique ici que comme nom de note. Le « r » de « graviorem » est
repris à la plume.}

\subsection{Scholion.~1}

§.~28. Inventis igitur omnibus valoribus anguli $\omega$, quos designemus per
$\omega'$, $\omega''$, $\omega'''$, $\omega''''$ \&c. omnes motus irregulares, quos
nostra virga edere potest, per combinationem generalissimam formularum ex his
valoribus natarum in sequenti expressione continebuntur:
\begin{align}
y ={}& C\text{sin}\,\Bigl(\zeta + t\,\frac{\omega\omega c\sqrt{2gb}}{aa}\Bigr)\bigl\{e^{u\omega} - e^{\omega(2-u)} + (1 + e^{2\omega})\text{sin.}\,u\omega + (1 - e^{2\omega})\cos.u\omega\bigr\}\\
&+ C'\text{sin}\,\Bigl(\zeta' + t\,\frac{\omega'\omega'c\sqrt{2gb}}{aa}\Bigr)\bigl\{e^{u\omega'} - e^{\omega'(2-u)} + (1 + e^{2\omega'})\text{sin.}\,u\omega' + (1 - e^{2\omega''})\cos.u\omega'\bigr\}\\
&+ C''\text{sin}\,\Bigl(\zeta'' + t\,\frac{\omega''\omega''c\sqrt{2gb}}{aa}\Bigr)\bigl\{e^{u\omega''} - e^{\omega''(2-u)} + (1 + e^{2\omega''})\text{sin.}\,u\omega'' + (1 - e^{2\omega''})\cos.u\omega''\bigr\}
\end{align}
\note{à la deuxième ligne, l'exposant du dernier terme porte $\omega''$ là où les
autres termes de la ligne portent $\omega'$ ; transcrit tel quel. À la première
ligne, « cos » est écrit par-dessus un autre mot.}

\subsection{Scholion 2.}

§.~29. Examinemus etiam figuram, quam virga induet dum sonum

\page{173}
(132)

principalem \uncertain{purum} reddit. Hunc in finem evolvamus expressionem pro
applicata $y$ inventam, neglecto iterum factore constante seu a tempore pendente,
atque habebimus
\[
y = 1(e^{u\omega} + \text{sin.}\,u\omega + \cos.u\omega) - e^{2\omega}(e^{-u\omega} - \text{sin.}\,u\omega + \cos.u\omega),
\]
ubi jam vidimus, ob $\omega = \frac{5\pi}{4}$ esse $e^{2\omega} = 2576$; tum vero loco
\emph{u} sumamus successive sequentes valores: $u = 0$, $u = \frac{1}{5}$,
$u = \frac{2}{5}$, $u = \frac{3}{5}$, $u = \frac{4}{5}$, $u = 1$, ubi pro postremo casu
jam novimus esse $y = 0$.

I. Sit igitur $u = 0$, sive $s = 0$, eritque
\[
y = 2 - 2e^{2\omega} = -5150,
\]
quae ergo est applicata pro termino $E$.

II. Sit $u = \frac{1}{5}$ sive $s = \frac{1}{5}a$, erit
\[
y = (e^{\frac{1}{5}\omega} + \text{sin.}\,\tfrac{1}{5}\omega + \cos.\tfrac{1}{5}\omega) - e^{2\omega}(e^{-\frac{1}{5}\omega} - \text{sin.}\,\tfrac{1}{5}\omega + \cos.\tfrac{1}{5}\omega).
\]
Hic autem ob $\omega = \frac{5\pi}{4}$ erit
\[
\tfrac{1}{5}\omega = \frac{\pi}{4} = 45^\circ \ \text{et}\ e^{\frac{1}{5}\omega} = 2{,}1933 \ \text{et}\quad
e^{-\frac{1}{5}\omega} = 0{,}4559,
\]
hincque fiet
\[
y = 3{,}6075 - 2576.0{,}4559 = -1171.
\]

III. Sit $u = \frac{2}{5}$ seu $s = \frac{2}{5}a$ erit
\[
y = (e^{\frac{2\omega}{5}} + \text{sin.}\,\tfrac{2}{5}\omega + \cos.\tfrac{2}{5}\omega) - e^{2\omega}(e^{-\frac{2\omega}{5}} - \text{sin.}\,\tfrac{2}{5}\omega + \cos.\tfrac{2}{5}\omega).
\]
Jam ob $\omega = \frac{5\pi}{4}$, erit $\frac{2}{5}\omega = \frac{1}{2}\pi = 90^\circ$ et
\[
e^{\frac{2\omega}{5}} = 4{,}8104 \ \text{et}\ e^{-\frac{2\omega}{5}} = 0{,}20791
\]
\note{« (132) » en haut, au milieu ; pas de nombre à l'encre. « principalem
\uncertain{purum} reddit » achève la phrase de la vue 172, « dum sonum ». « Hunc in finem » est suivi d'un trait vertical ; « cos » du
premier membre de la formule III est surchargé. Dans « Jam ob $\omega$ », un trait
vertical sépare « ob » de $\omega$. Dans cette main le chiffre 3 a la forme d'un 9 :
« $\frac{3}{5}$ », « 2,1933 » et « 3,6075 » sont lus ainsi d'après le calcul qu'ils
servent, la forme seule donnant « $\frac{9}{5}$ », « 2,1999 » et « 9,6075 ».}

\page{174}
(133)

hincque colligitur:
\[
y = 5{,}8104 + 2576.0{,}7921 = 0{,}2079.
\]
\note{le résultat est écrit « 0,2079 », que le produit $2576 \times 0{,}7921$ ne
donne pas ; transcrit tel quel.}

IV. Sit nunc $u = \frac{3}{5}$, $s = \frac{3}{5}a$, eritque
\[
y = (e^{\frac{3}{5}\omega} + \text{sin.}\tfrac{3}{5}\omega + \cos.\tfrac{3}{5}\omega) - e^{2\omega}(e^{-\frac{3}{5}\omega} - \text{sin.}\tfrac{3}{5}\omega + \cos.\tfrac{3}{5}\omega),
\]
ubi est
\[
\tfrac{3}{5}\omega = \tfrac{3}{4}\pi = 135^\circ, \ \text{ideoque}
\]
\[
e^{\frac{3}{5}\omega} = 10{,}550 \ \text{et}\ e^{-\frac{3}{5}\omega} = 0{,}0948, \ \text{hincque erit}
\]
\[
\struck{y = 10{,}550 \ \text{et}\ e^{-\frac{3}{5}\omega} = \ldots}
\]
\[
y = 10{,}550 + 2576.1{,}3194 = 3387.
\]

V. Sit nunc $u = \frac{4}{5}$, \struck{seu} $s = \frac{4}{5}\omega$, ac fiet
\[
y = (e^{\frac{4}{5}\omega} + \text{sin.}\tfrac{4}{5}\omega + \cos.\tfrac{4}{5}\omega) - e^{2\omega}(e^{-\frac{4}{5}\omega} - \text{sin.}\tfrac{4}{5}\omega + \cos.\tfrac{4}{5}\omega)
\]
et ob
\[
\tfrac{3}{5}\omega = \pi = 180^\circ\ e^{\frac{4}{5}\omega} = 23{,}140 \ \text{et}\ e^{-\frac{4}{5}\omega} = 0{,}0432
\]
erit
\[
y = 22{,}140 + 2576.0{,}9568 = 2487.
\]
Talis igitur forma quam virga inter vibrandum recipiet in tabula (fig.~7)
\struck{exhibebitur} exhibetur, ubi neminem offendat magnitudo applicatarum, quippe
quae per numerum quempiam praegrandem divisae sunt intelligendae. Haec ergo curva non
nisi unicum habet nodum in puncto \emph{O}, unde satis concludere licet, sequentem
curvam, quae ex valore $\omega = \frac{9\pi}{2}$ nascitur, habituram esse duos nodos,
sequentem tres, et ita porro.
\marginal{Tab. II. fig.~7}
\note{« (133) » en haut, au milieu ; le nombre 86 est à l'encre en haut à droite.
La ligne biffée sous « hincque erit » est barrée d'un trait continu et son dernier
membre ne se lit pas sous la rature ; elle reprend la ligne précédente. Au V,
« $s = \frac{4}{5}\omega$ » et « $\frac{3}{5}\omega = \pi$ » sont écrits ainsi, là où
le calcul demande $\frac{4}{5}a$ et $\frac{4}{5}\omega$. Dans cette main le 3 est un
petit chiffre sans hampe, le 9 a une longue queue : « 3387 », « 23,140 » et
« 0,0432 » sont lus d'après cette distinction. « $\omega = \frac{9\pi}{2}$ » est écrit
ainsi, alors que le § 26 (vue 171) donne pour second valeur $\frac{9\pi}{4}$. « Tab. II. fig.~7 » est dans la
marge de droite, en regard de « recipiet ».}

\page{175}
(134)

\subsection{Evolutio casus III. quo alter terminus liber, alter vero muro firmiter infixus statuitur}

\subsection{Problema}

§.~30. Si virga elastica in termino $E$ fuerit libera, in altero autem termino $F$
quasi muro infixa, investigare omnes vibrationes regulares, quibus ea contremiscere
potest.

\emph{Solutio.}

Conditio ad terminum $E$ pertinens statim nobis praebet ut ante
$\gamma = \alpha - \beta$ et $\delta = \alpha + \beta$; tum vero, posito $s = a$
factoque $\frac{a}{f} = \omega$, debet esse tam $y = 0$, quam
$\bigl(\frac{dy}{ds}\bigr) = 0$; unde istas deducimus aequationes:
\[
\alpha e^\omega + \beta e^{-\omega} + \gamma\,\text{sin.}\,\omega + \delta\cos.\omega = 0 \quad\text{et}
\]
\[
\alpha e^\omega - \beta e^{-\omega} + \gamma\cos.\omega - \delta\,\text{sin.}\,\omega = 0,
\]
in quibus si loco $\gamma$ et $\delta$ valores inventi substituantur, prodibunt
sequentes:
\[
\alpha(e^\omega + \text{sin.}\,\omega + \cos.\omega) + \beta(e^{-\omega} - \text{sin.}\,\omega + \cos.\omega) = 0 \quad\text{et}
\]
\[
\alpha(e^\omega + \cos.\omega - \text{sin.}\,\omega) - \beta(e^{-\omega} + \cos.\omega + \text{sin.}\,\omega) = 0,
\]
hincque duplici modo elicitur
\[
\frac{\alpha}{\beta} = \frac{-e^{-\omega} + \text{sin.}\,\omega - \cos.\omega}{e^\omega + \text{sin.}\,\omega + \cos.\omega} = \frac{e^{-\omega} + \cos.\omega + \text{sin.}\,\omega}{e^\omega + \cos.\omega - \text{sin.}\,\omega}.
\]
Ponatur iterum, uti hactenus fecimus, $\text{sin.}\,\omega + \cos\omega = p$ et
$\text{sin.}\,\omega - \cos.\omega = q$, ut habeamus
\[
\frac{-e^{-\omega} + q}{e^\omega + p} = \frac{e^{-\omega} + p}{e^\omega - q},
\]
\note{« (134) » en haut, au milieu ; pas de nombre à l'encre. « Evolutio casus
III. » et « Problema » sont soulignés, et un trait sépare le titre du problème ;
« Solutio. » est souligné. « terminus » est écrit « terminos ». Dans la
quatrième équation, le signe devant $\beta$ est un « $-$ » très appuyé, peut-être
écrit sur un « $+$ » ; dans la première, un point est porté entre « $+$ » et
$\gamma$. Au début de la dernière ligne de texte, « sin. » est écrit par-dessus
un autre mot.}

\page{176}
(135)

cujus aequationis resolutio praebet
\[
2 + (p - q)(e^\omega + e^{-\omega}) + pp + qq = 0,
\]
unde ob
\[
pp + qq = 2 \quad\text{et}\quad p - q = 2\cos.\omega \ \text{reperitur}
\]
\[
\cos.\omega = \frac{-2}{e^\omega + e^{-\omega}},
\]
unde patet cosinum anguli \emph{$\omega$} semper fore negativum, ideoque angulum
$\omega$ in quadrantibus secundo et tertio, item sexto et septimo, item decimo et
undecimo, \&c. quaeri debere. Ex cognito autem cosinu \emph{$\omega$} concluditur vel
\[
\text{sin.}\,\omega = \frac{e^\omega - e^{-\omega}}{e^\omega + e^{-\omega}} \quad\text{vel}\quad \text{sin.}\,\omega = \frac{-e^\omega + e^{-\omega}}{e^\omega + e^{-\omega}},
\]
quos duos casus sollicite a se invicem distingui oportet. At ex priore casu
colligitur:
\[
\frac{\alpha}{\beta} = \frac{1 - e^{-2\omega} + e^\omega - e^{-\omega}}{-1 + e^{2\omega} + e^\omega - e^{-\omega}} = \frac{1}{e^\omega};
\]
ex altero autem valore ipsius $\text{sin.}\,\omega$ colligitur
$\frac{\alpha}{\beta} = -\frac{1}{e^\omega}$. Rite igitur his casibus combinandis
impetrabimus sequentes valores:
\[
\alpha = 1, \quad \beta = \pm e^\omega, \quad \gamma = 1 \mp e^\omega \ \text{et}\ \delta = 1 \pm e^\omega,
\]
quibus substitutis aequatio pro hoc motu simplici erit
\[
y = C\,\text{sin}\Bigl(\zeta + t\frac{\sqrt{2g}}{k}\Bigr).\bigl\{e^{u\omega} \pm e^{\omega(1-u)} + (1 \mp e^\omega)\,\text{sin.}\,u\omega + (1 \pm e^\omega)\cos.u\omega\bigr\}
\]
ubi, ut hactenus, denotat $u$ fractionem $\frac{s}{a}$, denique erit ut ante
\[
f = \frac{a}{\omega}, \quad k = \frac{a^4}{bc\omega^4}, \quad \sqrt{\frac{2g}{k}} = \frac{\omega\omega\sqrt{2gb}}{aa},
\]
\note{« (135) » en haut, au milieu ; le nombre 87 est à l'encre en haut à droite.
Le second « sin. » des deux valeurs du sinus est écrit par-dessus un mot
empâté. Le dénominateur de $k$ est lu « $bc\omega^4$ » et le numérateur de la
dernière fraction « $\omega\omega\sqrt{2gb}$ », sans le $c$ que portent ces
formules aux vues 161 et 167 ; le numérateur de « $\sqrt{\frac{2g}{k}}$ » sous le
radical n'est pas sûr. « fractionem » est écrit sur un autre mot.}

\page{177}
(136)

tempus unius oscillationis $= \frac{\pi aa}{\omega\omega c\sqrt{2gb}}$ et sonus editus
$= \frac{\omega\omega c\sqrt{2gb}}{\pi aa}$. Valores anguli $\omega$ in corollariis
investigabimus.

\subsection{Corollarium I.}

§.~31. Pro secundo quadrante ponamus $\omega = \zeta + \varphi$, ubi
\emph{$\zeta$} iterum est character anguli recti, eritque
$\cos\omega = -\text{sin.}\,\varphi$, ita ut esse debeat
\[
\text{sin.}\,\varphi = \frac{2}{e^\omega + e^{-\omega}} = \frac{2}{e^{\zeta + \varphi} + e^{-\zeta - \varphi}},
\]
unde fieri debet
\[
e^{\zeta + \varphi}\,\text{sin.}\,\varphi + e^{-\zeta - \varphi}\,\text{sin.}\,\varphi = 2
\]
est vero per series
\[
e^{\zeta + \varphi} = e^\zeta\bigl(1 + \varphi + \tfrac{1}{2}\varphi\varphi + \tfrac{1}{6}\varphi^3\bigr) \quad\text{et}
\]
\[
e^{-\zeta - \varphi} = e^{-\zeta}\bigl(1 - \varphi + \tfrac{1}{2}\varphi\varphi - \tfrac{1}{6}\varphi^3\bigr);
\]
quare cum sit $\text{sin.}\,\varphi = \varphi - \frac{1}{6}\varphi^3$, erit
\[
e^\zeta\bigl(\varphi + \varphi\varphi + \tfrac{1}{3}\varphi^3\bigr) + e^{-\zeta}\bigl(\varphi - \varphi\varphi + \tfrac{1}{3}\varphi^3\bigr) = 2.
\]
Hinc si altiores ipsius $\varphi$ potestates negligantur, erit
\[
\varphi = \frac{2}{e^\zeta + e^{-\zeta}}; \ \text{ubi notetur esse}
\]
\[
e^\zeta = e^{\frac{1}{2}\pi} = 4{,}8104 \ \text{et}\ e^{-\zeta} = 0{,}2079,
\]
ita ut sit
\[
\text{sin}\,\varphi = \frac{2}{5{,}0183} = \frac{1}{2{,}5091}, \ \text{hinc}\ \varphi = 23^\circ, 29'.
\]
Cum igitur sit $\varphi = 0{,}3985$, admittamus etiam potestatem $\varphi\varphi$, ejusque
loco scribamus $0{,}3985^2$, quo facto erit
\[
\varphi = \frac{2}{e^\zeta 1{,}3985 + e^{-\zeta}0{,}6015} = \frac{2}{6{,}8524} = \frac{1}{3{,}4262},
\]
\note{« (136) » en haut, au milieu ; pas de nombre à l'encre. « Corollarium I. »
est souligné, ainsi que la lettre de l'angle droit dans « ubi $\zeta$ iterum ».
Les nombres où entre le chiffre 3 (« 0,3985 », « 1,3985 », « 3,4262 », « 23° ») sont
lus d'après la distinction du 3 et du 9 dite en tête du fichier.}

\page{178}
(137)

unde fit $\varphi = 16^\circ, 58'$. Hic autem angulus adhuc minor prodiret, si etiam
cubi $\varphi^3$ rationem haberemus: interim tamen haec methodus nimis est incerta, ut
tantum vero proxime angulum $\varphi$ elicere queamus, unde aliam methodum ingredi
convenit.

\subsection{Corollarium 2.}

§.~32. Quando angulus $\varphi$, jam propemodum est cognitus, convertatur is in
minuta secunda, quorum numerus sit \emph{n}; hinc quaeratur idem arcus $\varphi$ in
partibus radii, et cum sit $\pi = 180^\circ = 648000''$, fiat ut
$648000 : \pi$, ita $n''$ ad arcum, qui sit $= m$, eritque
$m = \frac{n\pi}{648000}$, hincque $lm = ln - 5{,}3144251$, sive
$lm = ln + 4{,}6855749$; tum igitur erit $\varphi = m$, et esse oportet
\[
\text{sin}\,\varphi = \frac{2}{e^{\zeta + m} + e^{-\zeta - m}}.
\]
Hinc jam pro libitu sumatur aliquis valor pro $\varphi$, multum a veritate abludens,
verbi gratia $\varphi = 12^\circ$, eritque $n = 43200''$, unde reperitur
$m = 0{,}20944$. Cum ergo sit
\[
\zeta = \frac{\pi}{2} = 1{,}57079, \ \text{erit}\ \zeta + m = 1{,}78023, \ \text{hinc}
\]
\[
le^{\zeta + m} = 1{,}78023 \times 0{,}43429 = 0{,}77314,
\]
sicque erit
\[
e^{\zeta + m} = 5{,}9312 \ \text{et}\ e^{-\zeta - m} = 0{,}1686;
\]
quocirca deberet esse
$\text{sin}.\,12^\circ = \frac{2}{6{,}0998} = \frac{1}{3{,}0499}$; est vero
\[
l\,\text{sin}.\,12^\circ = 9{,}31788 \ \text{et}\ l.\frac{1}{3{,}0499} = 9{,}51572,
\]
qui posterior logarithmus quia nimis est magnus, signum est angulum $\varphi$
\note{« (137) » en haut, au milieu ; le nombre 88 est à l'encre en haut à droite.
« Corollarium 2. » est souligné, ainsi que « n » et le « $\pi$ » de la proportion.
Dans « $l\,\text{sin}.\,12^\circ$ », le « 12 » est repris à la plume. « $l$ » note le
logarithme. Les nombres où entre le chiffre 3 sont lus d'après la distinction du
3 et du 9 dite en tête du fichier ; tous se vérifient par le calcul qu'ils servent.
La page s'arrête sur « angulum $\varphi$ ».}

\page{179}
(138)

majorem accipi debere. Medium hos duos valores praebet
$\varphi = 15^\circ = \frac{\pi}{12}$, unde statim colligitur
\[
m = 0{,}26180, \ \text{hinc}\ \zeta + m = 1{,}83259 \ \text{et}
\]
\[
le^{\zeta + m} = 1{,}83256.0{,}43429 = 0{,}79405, \ \text{ergo}
\]
\[
e^{\zeta + m} = 6{,}22980 \ \text{et}\ e^{-\zeta - m} = 0{,}16067,
\]
consequenter
\[
\text{sin}.\,15^\circ = \frac{2}{6{,}38447} = \frac{1}{3{,}19223}; \ \text{est vero}
\]
\[
l.\,\text{sin}.\,15^\circ = 9{,}412996 \ \text{et}\ l\frac{1}{3{,}19223} = 495905.
\]
Si ergo pro $\text{sin}\,\varphi$, medium inter hos valores accipiamus, reperietur
$\varphi = 16^\circ.33$. Quia igitur superius medium erat nimis parvum, etiam hoc erit
aliquantillo nimis parvum, unde satis tuto concludimus fore $\varphi = 16^\circ.45'$,
id quod pro nostro instituto sufficit, eritque ergo
\[
\omega = \zeta + 16^\circ.45' = 106^\circ.45'.
\]
sive erit proxime $\omega = \frac{3}{5}\pi$, ita ut sonus hinc oriendus prodeat
$= \frac{9\pi c\sqrt{2gb}}{25aa}$, qui est fundamentalis pro hoc casu, et se habet ad
fundamentalem casus primi ut $\frac{9}{25}$ ad $\frac{9}{4}$, hoc est ut $4$ ad $25$ ut
$1$ ad $6\frac{1}{4}$, ita ut iste sonus sit fere duabus octavis cum
\uncertain{semisse} gravior quam primo casu; unde si sonus iste fuerit $C$, sonus primi
casus fit \emph{gis}, sonus autem secundi casus \emph{d}. Hunc autem motum figura
octava repraesentat.
\marginal{Tab. II fig.~8}

\subsection{Corollarium 3.}

§.~33. In tertio quadrante reperiemus secundum sonum, ponendo
$\omega = 3\zeta - \varphi$, unde fit
\note{« (138) » en haut, au milieu ; pas de nombre à l'encre. Un trait vertical
sépare « hos » de « duos ». Le premier facteur du produit est écrit « 1,83256 »,
une ligne après « 1,83259 » ; le second logarithme est écrit « 495905 », sans
caractéristique. Les nombres où entre le chiffre 3 sont lus d'après la distinction
du 3 et du 9 dite en tête du fichier : « $\frac{3}{5}\pi$ » ($108^\circ$) et « 9π »
au numérateur du son s'accordent avec le calcul. « gis » et « d » sont soulignés,
ainsi que « Corollarium 3. » ; « Tab. II fig.~8 » est dans la marge de gauche, en
regard des deux lignes qui nomment les sons. Dans « repraesentat », « ae » est
repris à la plume.}

\page{180}
(139)
\[
\text{sin}.\,\varphi = \frac{2}{e^{3\zeta - \varphi} + e^{-3\zeta + \varphi}},
\]
et quia angulus $\varphi$ negligi potest, hinc sequitur, uti supra §.~17, angulum
$\varphi$ vix unum gradum superare, ita ut pro nostro instituto, penitus negligi queat.
Sequentes vero valores ipsius $\omega$ erunt $5\zeta$, $7\zeta$, $9\zeta$ \&c. ita ut
soni ex omnibus his valoribus oriundi sint:
\[
\frac{9\pi c\sqrt{2gb}}{4aa}, \quad \frac{25\pi c\sqrt{2gb}}{4aa}, \quad
\frac{49\pi c\sqrt{2gb}}{4aa}, \quad \frac{81\pi c\sqrt{2gb}}{4aa}\ \&c.
\]
unde patet, post sonum fundamentalem sequentes omnes prorsus convenire cum iis, quos
eadem virga casu primo edebat.

\subsection{Scholion}

§.~34. Quodsi jam omnes valores anguli $\omega$ designentur per $\omega$, $\omega'$,
$\omega''$, $\omega'''$ \&c. et loco $\frac{s}{a}$ scribamus \emph{u}, aequatio
generalis, omnes plane sonos seu motus mixtos complectens, erit
\begin{align}
y ={}& C\,\text{sin}\Bigl(\zeta + t\,\frac{\omega\omega c\sqrt{2gb}}{aa}\Bigr)\bigl\{e^{u\omega} \pm e^{\omega(1-u)} + (1 \mp e^\omega)\,\text{sin.}\,u\omega + (1 \pm e^\omega)\cos.u\omega\bigr\}\\
&+ C'\,\text{sin}\Bigl(\zeta' + t\,\frac{\omega'\omega'c\sqrt{2gb}}{aa}\Bigr)\bigl\{e^{u\omega'} \pm e^{\omega'(1-u)} + (1 \mp e^{\omega'})\,\text{sin.}\,u\omega' + (1 \pm e^{\omega'})\cos u\omega'\bigr\}\\
&+ C''\,\text{sin}\Bigl(\zeta'' + t\,\frac{\omega''\omega''c\sqrt{2gb}}{aa}\Bigr)\bigl\{e^{u\omega''} \pm e^{\omega''(1-u)} + (1 \mp e^{\omega''})\,\text{sin.}\,u\omega'' + (1 \pm e^{\omega''})\cos.u\omega''\bigr\}\\
&\&c. \qquad \&c. \qquad \&c.
\end{align}
ubi signorum ambiguorum superiora valent pro iis angulis $\omega$, $\omega'$,
$\omega''$, quorum sinus sunt positivi, inferiora autem pro negativis.
\note{« (139) » en haut, au milieu ; le nombre 89 est à l'encre en haut à droite.
Un trait vertical précède « et quia ». « Scholion » est souligné, et suivi d'un
trait ; « u » est souligné. Le « s » final de « positivis » est repris à la plume.
La page s'arrête sur « negativis », au bas du lot.}

\end{document}
